\section{Introduction}
\label{sec:intro}

Federated learning (FL) was conceived as a way to learn from data that
cannot be moved: hospitals coordinate without sharing
records~\cite{sheller2020medicine,dayan2021covid}, pharmaceutical
companies coordinate without disclosing proprietary chemistry~\cite{melloddy2024},
and industrial-scale deployments coordinate across millions of
devices~\cite{bonawitz2019sysml}. The protective premise is simple:
clients share model updates, not raw data, so privacy follows from the
information bottleneck of communication. The literature on privacy
\emph{within} that bottleneck is rich. Gradient inversion can reconstruct
training samples from un-noised gradient observations~\cite{zhu2019dlg,geiping2020inverting};
membership inference reveals participation given black-box query
access~\cite{shokri2017membership}; collaborative-learning leakage
extracts data properties from the joint update
sequence~\cite{melis2019exploiting}. Differentially private stochastic
gradient descent (DP-SGD)~\cite{abadi2016dpsgd} and its federated variant
DP-FedAvg~\cite{mcmahan2018dpfedavg} bound this content-level leakage and
have become the canonical defense, with R\'enyi-DP composition giving
the tightest practical accounting~\cite{mironov2017rdp}.

What that programme does \emph{not} address is the topology around the
content. Real deployments organise clients into structured graphs, including
hierarchical client--edge--cloud cascades~\cite{liu2020hierarchical,abad2020icassp},
gossip-based peer-to-peer protocols~\cite{roy2019braintorrent,hegedus2021decentralized},
and DAG-aggregation variants~\cite{beilharz2021dag}, because no single
flat protocol meets every constraint. The topology choice is treated
as a performance lever (bandwidth, latency, scalability)~\cite{wang2021fieldguide};
its privacy implications have remained largely unexamined. An adversary
who observes per-round updates and knows the topology exploits two
channels that DP-SGD does not bound: the structural position of each
client in the communication graph, and the organisational membership
encoded by deployment design.

This paper addresses that gap with two contributions, an attack and a
defense, each validated within the same threat model. We identify a
class of \emph{topology-conditional} distributional inference attacks
against DP-protected FL, prove a per-client mutual-information bound
that applies against these adversaries, and derive a noise allocation
that is provably optimal under that bound.

\paragraph{Attack: \TADI.}
\TADI (Topology-Aware Distributional Inference) is a
shadow-trained~\cite{shokri2017membership} regressor that maps
\emph{(parameter sequence, structural feature vector, organisational
label)} to an estimate of each client's sensitive-class concentration
$p_i$. We instantiate four channel ablations
($\Adv_1, \Adv_2^{\text{topo}}, \Adv_2^{\text{org}}, \Adv_2^{\text{full}}$)
that isolate the contribution of each information channel, and we
calibrate against a constant-mean baseline so that the IID-null
condition (no topology--data correlation) gives zero attack lift by
construction.

\paragraph{Defense: \Fulcrum.}
The defense takes its name from the lever's balance point: under a
fixed utility budget, noise is allocated asymmetrically so that
per-client privacy bounds are equilibrated across the leverage profile,
with more noise assigned to high-leverage clients and less to
low-leverage ones. The key structural insight is that the per-client
conditional mutual information splits into two additive terms, one
shrinkable by DP noise and one not, and the optimal allocation targets
only the shrinkable term:
\begin{equation*}
I(p_i;\hat{p}_i \mid \TOPO,\omega,\{\sigma_j\})
\;\leq\;
\underbrace{\frac{\Tmax C^2}{2\sigma_i^2 |B|^2}}_{\text{controllable}}
\;+\;
\underbrace{\ell_i^\circ}_{\text{prior coupling}}
\end{equation*}
By Theorems~\ref{thm:per-client} and~\ref{thm:allocation}, the
closed-form min-max optimal allocation
$\sigma_i^{*2} = a/(K^\star - \ell_i^\circ)$ \emph{strictly dominates}
the uniform DP-SGD allocation whenever the structural-leverage scores
$\ell_i^\circ$ are non-uniform across clients. Three practical proxies
cover the canonical deployment regimes.

\paragraph{Contributions.}
\begin{enumerate}[leftmargin=*,topsep=2pt,itemsep=2pt]
\item A formal threat model for topology-conditional distributional
  inference under DP-SGD, with calibration loss and attack lift over a
  constant-mean baseline as the primary per-client metrics, designed
  to retain statistical power at small client counts
  (\S\ref{sec:threat}).
\item The \TADI attack with four channel ablations and a constant-mean
  comparator, providing the first per-channel decomposition of
  topology-conditional leakage in DP-FL (\S\ref{sec:attack}).
\item Theorems~\ref{thm:per-client} and~\ref{thm:allocation}
  establishing an additive MI bound and a closed-form balanced min-max
  allocation, with three principled leverage proxies and an explicit
  regime-of-applicability characterisation (\S\ref{sec:defense}).
\item Empirical evaluation on Fed-ISIC2019, Fed-Heart-Disease, and
  synthetic CIFAR-10 with parametric coupling, validating the defense
  end-to-end (\S\ref{sec:experiments}). Topology-aware allocation
  strictly dominates uniform DP-SGD on the privacy bound at every
  $(U, \Tmax)$ cell tested; the defense costs no measurable utility
  (paired $t$-tests over seeds give $p > 0.48$ at every cell); and
  DP-SGD bounds the \TADI parameter channel ($\Adv_1$ lift negative)
  across all settings. The prior-coupling term of
  Theorem~\ref{thm:per-client} is empirically characterised across
  deployment regimes: when the adversary's shadow prior matches the
  target deployment's prior (Setting~C, $\eta$-coupled synthetic), the
  bound is attained, with attack lift positive, monotone in the
  coupling parameter, and calibration-null at $\eta=0$. Under the
  public-proxy threat model on cross-silo benchmarks (Settings~B
  and~A), where the adversary has only synthetic re-partitioning of
  public data, the bound is conservative in a deployment-favourable
  direction (\S\ref{sec:discussion}).
\end{enumerate}

\paragraph{Out of scope.}
We assume an honest-but-curious passive adversary; active manipulation
of the protocol, malicious aggregators, and Sybil attacks are
complementary threats outside our scope.
Secure aggregation~\cite{bonawitz2017secagg} reduces our adversary's
observations to neighbourhood sums and changes the threat model;
we evaluate this as an extension experiment rather than the headline
threat. Our results assume static topologies and disjoint client
datasets; dynamic reconfiguration and overlapping client populations
are flagged as limitations in \S\ref{sec:discussion}.
